% ============================================================================
% Section 3: Pipeline Architecture
% ============================================================================
\section{The RAMA Discovery Pipeline}
\label{sec:pipeline}

\subsection{Overview}

The \RAMA{} (\emph{Ramanujan Autonomous Neuro-symbolic
Architecture}) pipeline consists of five sequential stages,
each implemented as a self-contained module with explicit
inputs, outputs, and failure modes.

\begin{enumerate}[label=\textbf{Stage~\arabic*:},leftmargin=3cm]
  \item \textbf{Vision Extraction.}
    A multimodal vision model (Gemini~2.5 Flash) processes
    each digitized manuscript page and returns a structured
    Pydantic schema containing: raw \LaTeX{} transcription,
    archetype classification (mock theta function,
    continued fraction, hypergeometric series, or
    $q$-series), confidence score, and---when
    identifiable---the first 12 $q$-series coefficients.

  \item \textbf{Genetic Evolutionary Search.}
    A population-based genetic algorithm searches the
    combinatorial space of eta-quotient exponent vectors
    $\mathbf{r} = (r_1, r_2, \ldots, r_{d_{\max}})$ with
    $d_{\max} = 12$ and $r_d \in [-24, 24]$.

  \item \textbf{Deep Think Bridge.}
    The archetype classification from Stage~1 is used to
    assign a target domain (String Theory / K3, Topology,
    or Astrophysics) and a conjectured modular shadow.

  \item \textbf{Lean~4 Auto-Formalization.}
    The code generator produces a Lean~4 file containing:
    \begin{itemize}
      \item \emph{Tier~A}: an \texttt{EtaQuotient} structure
            with exact rational arithmetic definitions for
            $\ceff$, weight, and leading power, verified by
            \texttt{rfl};
      \item \emph{Tier~B}: a structural blueprint encoding
            the shadow and domain classification as typed
            data.
    \end{itemize}
    The Lean verifier invokes \texttt{lake env lean} on
    the generated file.  Compilation with exit code~0
    constitutes \emph{kernel certification}.

  \item \textbf{Physical Mapping.}
    Verified candidates are evaluated via a saddle-point
    matrix computing BPS state entropy, effective central
    charge, and thermodynamic stability classification.
\end{enumerate}

\subsection{Genetic Algorithm Details}

\subsubsection{State Representation}

Each individual in the population is an
\texttt{EtaQuotientState}:
\begin{equation}
  \mathcal{S} = \bigl(
    (r_d)_{d=1}^{d_{\max}},\;
    q_{\mathrm{shift}} \in [-72, 72]
  \bigr),
\end{equation}
where $r_d \in \{-24, -23, \ldots, 23, 24\}$ and
$q_{\mathrm{shift}}$ is the fractional power shift
in units of $1/24$.

\subsubsection{Fitness Function}

The composite energy functional is
\begin{equation}
  \label{eq:energy}
  E(\mathcal{S}) = \alpha \cdot C(\mathcal{S})
  + \beta \cdot I(\mathcal{S})
  + \gamma \cdot D(\mathcal{S}),
\end{equation}
where:
\begin{itemize}
  \item $C(\mathcal{S})$ is the \emph{complexity} cost:
        a monotonic function of $\sum |r_d|$ penalizing
        overly complex quotients (Occam regularization);
  \item $I(\mathcal{S})$ is the \emph{fit error}: the
        normalized $\ell_2$ distance between the first
        $n$ coefficients of the candidate $q$-expansion
        and the target coefficients extracted by Stage~1;
  \item $D(\mathcal{S})$ is the \emph{domain penalty}:
        a term encoding prior knowledge about physically
        viable exponent ranges.
\end{itemize}

\subsubsection{Evolutionary Operators}

\begin{itemize}
  \item \textbf{Selection}: Tournament selection with $k=3$.
  \item \textbf{Crossover}: Uniform crossover of exponent vectors.
  \item \textbf{Mutation}: Each $r_d$ perturbed with probability
        $p_{\mathrm{mut}} = 0.3$ by $\Delta r_d \sim \mathrm{Uniform}(-4, 4)$.
  \item \textbf{Immigration}: When diversity drops below
        $|\{\text{unique configs}\}| < 5$, five random
        immigrants are injected.
\end{itemize}

\subsubsection{Hyperparameters}

For the full manuscript run:
\begin{center}
\begin{tabular}{lcc}
\toprule
\textbf{Parameter} & \textbf{10-page dry run} & \textbf{Full run} \\
\midrule
Population size & 25 & 50 \\
Max generations & 5 & 15 \\
$d_{\max}$ & 12 & 12 \\
Mutation rate & 0.3 & 0.3 \\
Tournament size $k$ & 3 & 3 \\
Lean retries & 2 & 2 \\
\bottomrule
\end{tabular}
\end{center}

\subsection{Lean~4 Verification Protocol}

The Lean~4 verification constitutes the formal gate of the
pipeline.  We use \textbf{Lean 4} with \textbf{Mathlib4}
(\texttt{leanprover/lean4:leanprover/lean4:v4.x})~\cite{moura_lean4,lean4_mathlib}.

\begin{definition}[Kernel Certification]
\label{def:kernel_cert}
A candidate eta-quotient is \emph{kernel-certified} if:
\begin{enumerate}
  \item The generated Lean~4 file compiles with
        \texttt{lake env lean} returning exit code~0.
  \item The file contains no \texttt{sorry}, \texttt{axiom},
        or \texttt{native\_decide} calls.
  \item The CI audit script
        (\texttt{dualscale/ci/audit\_lean.py}) reports
        no forbidden axioms.
\end{enumerate}
\end{definition}

\begin{remark}
Kernel certification confirms \emph{type-level correctness}
of the stated algebraic identities.  It does \emph{not}
verify that the candidate eta-quotient matches Ramanujan's
original intent or that it is ``new'' relative to the
existing literature.  Such claims require separate
verification against databases such as the
OEIS~\cite{ono2004} and LMFDB.
\end{remark}
